\documentclass[12pt, letterpaper, parskip = full]{article}
\usepackage[mmddyyyy]{datetime}

\usepackage[utf8]{inputenc}
\usepackage{amsmath}
\usepackage{mathpazo}
\usepackage{eulervm}
\usepackage[skip=10pt plus1pt]{parskip}
\usepackage[normalem]{ulem}

\usepackage[margin=1in]{geometry}
\linespread{1.20}
\usepackage{enumerate}
\usepackage{enumitem}
\usepackage{listings}
\usepackage[T1]{fontenc}
\usepackage{amsfonts}
\usepackage{amssymb}
\usepackage{amsthm}
\usepackage{multirow,array}
\usepackage{graphicx}
\usepackage{physics}
\usepackage{bm}
\usepackage{tikz}
\usepackage{breqn}
\usepackage[longnamesfirst]{natbib}
\usepackage{comment}
\usepackage{setspace}
\usepackage{hyperref}
\usepackage{accents}
\usepackage{graphicx}

\graphicspath{{Images/}}
\setlength{\parindent}{0pt}

\newtheorem{theorem}{Theorem}
\newtheorem{corollary}{Corollary}
\newtheorem{lemma}{Lemma}
\newtheorem{definition}{Definition}
\newtheorem{axiom}{Axiom}
\newcommand{\e}{\mathrm{e}}
\newcommand{\Var}{\mathrm{Var}}
\newcommand{\Ex}[1]{\mathbb{E}\left[#1\right]}
\renewcommand{\Pr}[1]{\mathbb{P}(#1)}
\newcommand{\plim}{\mathrm{plim}}
\newcommand{\pc}[1]{\frac{\dot{#1}(t)}{#1(t)}}
\renewcommand\qedsymbol{$\blacksquare$}
\newcommand{\ubar}[1]{\underaccent{\bar}{#1}}

\title{Worksheet 8}
\author{Alejandro Flores-Brown}

\begin{document}

	\maketitle

    \section{W 8-1}

    The initial cashflow is \(\$100000\), the project's life is \(5\) years, we expect \(100,000\) cupcakes to be sold per year with an expected price of \(\$0.75\) per unit sold and an expected variable cost of \(\$0.25\) per unit sold. Annual fixed costs are \(\$10,000\) and taxes are \(30\%\). 

    \begin{center}
	\begin{tabular}{|c|c|c|}
		\hline
		Year & 0 & 1-5 \\
		\hline
		Sales&  & \(75000\) \\
        \hline
        - Costs & & \(-35000\) \\
        \hline
        - Depreciation & & \(-20000\) \\
        \hline
        \hline
        EBIT & & \(20000\) \\
        \hline
        - taxes & & \(-6000\) \\
        \hline
        \hline
        Net Income & & \(14000\) \\
        \hline
        + Depreciation & & \(20000\) \\
        \hline
        \hline
        Operating CF & & \(34000\) \\
        \hline
        - capital spending & \(-100000\) & \(0\) \\
        \hline
        - change in NWC & & \(0\) \\
        \hline
        \hline
        Project CF & \(-100000\) & \(34000\) \\
        \hline
	\end{tabular}
	\end{center}

    NPV: 
    \[ \Ex{R} = -100000 + \frac{34000}{1.1} + \frac{34000}{1.1^2} + \frac{34000}{1.1^3} + \frac{34000}{1.1^4} + \frac{34000}{1.1^5} = 28886.75 \]

    \section{W 8-2}

    \subsection{part A}

    We now want to conduct a sensitivity analysis. Let the units sold be \(u \in (70000,100000,130000)\). Build a table.

    \begin{center}
	\begin{tabular}{|c|c|c|c|c|}
		\hline
		Year & 0 & High Sales & Medium Sales (base case) & Low Sales \\
		\hline
		Sales & & \(130000\) & \(75000\) & \(52500\) \\
        \hline
        - Costs & & \(-42500\) & \(-35000\) & \(-27500\)\\
        \hline
        - Depreciation & & \(-20000\) &  \(-20000\) & \(-20000\) \\
        \hline
        \hline
        EBIT & & \(67500\) & \(20000\) & \(5000\) \\
        \hline
        - taxes & & \(-20250\) & \(-6000\) & \(-1500\) \\
        \hline
        \hline
        Net Income & & \(47250\) & \(14000\) & \(3500\) \\
        \hline
        + Depreciation & & \(20000\) & \(20000\) & \(20000\) \\
        \hline
        \hline
        Operating CF & & \(67250\) & \(34000\) & \(23500\) \\
        \hline
        - capital spending & \(-100000\) & \(0\) & \(0\) & \(0\) \\
        \hline
        - change in NWC & & \(0\) & \(0\) & \(0\) \\
        \hline
        \hline
        Project CF & \(-100000\) & \(67250\) & \(34000\) & \(23500\) \\
        \hline
        \hline
        NPV & & \(154930.41\) & \(28886.75\) & \(-10916.51\) \\
        \hline 
	\end{tabular}
	\end{center}

    We could also back-solve these equations to find how many sales we need to break even. This comes out to \(75485\) sales per year.

    \subsection{part B}

    We now want to do a scenario analysis. The idea behind a scenario analysis is to look at several changes in economic conditions to see how your cashflows might change. (e.g. scenario 1 here has high demand, higher per unit price, and lower variable costs). [ran out of time in class for this problem]

\end{document}